\subsection*{Цели и задачи}
\textit{Цель работы:} изучение характера и степени искажений импульсных сигналов прямоугольной формы при прохождении их через линейные цепи.

\textit{Задачи:}
\begin{itemize}
    \item Исследовать искажения в цепях первого и второго порядков, применяя временной подход (оценка через длительность переходного процесса).
    \item Исследовать искажения в цепях высокого порядка, применяя частотный подход (оценка через ширину спектра и полосу пропускания).
\end{itemize}

\subsection*{Краткие теоретические сведения}
Искажения прямоугольных импульсов проявляются в удлинении фронтов, появлении колебаний (выбросов) и спаде плоской вершины.

Для цепи первого порядка собственная частота определяется выражением:
\begin{equation}
    p_1 = -1/(RC) = -1/\tau_{\text{ц}}
\end{equation}
где $\tau_{\text{ц}}$ --- постоянная времени цепи.

Для цепи второго порядка собственные частоты вычисляются как:
\begin{equation}
    p_{1,2} = -\frac{1}{2RC} \pm \sqrt{\left(\frac{1}{2RC}\right)^2 - \frac{1}{LC}}
\end{equation}

В цепях высокого порядка используется частотный подход. Граничная частота (частота среза) $f_{\text{ср}} = \omega_{\text{ср}} / 2\pi$ определяется из условия:
\begin{equation}
    |H_U(j\omega_{\text{ср}})| = H_{U\max} / \sqrt{2}
\end{equation}

Ширина спектра прямоугольного импульса длительностью $t_{\text{и}}$ по первому нулю спектра:
\begin{equation}
    \Delta\omega_{\text{сп}} = 2\pi / t_{\text{и}}; \quad \Delta f_{\text{сп}} = 1 / t_{\text{и}}
\end{equation}

Искажения будут малы при выполнении условия:
\begin{equation}
    \Delta f_{\text{сп}} < f_{\text{ср}}
\end{equation}

\subsection*{Инструкция по снятию данных}

\textit{1. Исследование искажений в цепи первого порядка}
\begin{enumerate}
    \item Соберите схему цепи первого порядка. Для этого используйте элементы с параметрами: $C = 200~\text{пФ}$, $R = 5~\text{кОм}$.
    \item Настройте генератор сигналов. Установите генерацию прямоугольных импульсов.
    \item Настройте осциллограф. Установите ждущий режим (Norm или Sing).
    \item Подайте на вход сигнал длительностью $t_{\text{и}} = 2~\text{мкс}$. Зафиксируйте осциллограммы входного $u_{\text{вх}}(t)$ и выходного $u_{\text{вых}}(t)$ напряжений. Сохраните данные.
    \item Измените длительность входного сигнала на $t_{\text{и}} = 10~\text{мкс}$. Повторно снимите и сохраните осциллограммы.
    \item \textit{Нюанс:} Убедитесь, что масштабы времени и амплитуды на осциллографе позволяют четко рассмотреть фронты сигналов.
\end{enumerate}

\begin{figure}[H]
    \centering
    \includegraphics[width=0.6\textwidth]{figs/circuit_1st_order.pdf}
    \caption{Схема цепи первого порядка}
    \label{fig:circuit_1st_order}
\end{figure}

\textit{2. Исследование искажений в цепи второго порядка}
\begin{enumerate}
    \item Соберите схему цепи второго порядка. Базовые параметры реактивных элементов: $L = 360~\text{мкГн}$, $C = 200~\text{пФ}$.
    \item Установите длительность входного прямоугольного импульса $t_{\text{и}} = 10~\text{мкс}$.
    \item Установите сопротивление $R_1 = 4~\text{кОм}$. Снимите осциллограммы $u_{\text{вх}}(t)$ и $u_{\text{вых}}(t)$.
    \item Замените сопротивление на $R_2 = 0,67~\text{кОм}$. Повторите снятие осциллограмм.
    \item Замените сопротивление на $R_3 = 0,1~\text{кОм}$. Повторите снятие осциллограмм.
    \item \textit{Предупреждение:} В данном пункте ожидается наблюдение различных режимов переходного процесса (апериодический, критический, колебательный). Тщательно фиксируйте наличие выбросов.
\end{enumerate}

\begin{figure}[H]
    \centering
    \includegraphics[width=0.4\textwidth]{figs/circuit_2nd_order.pdf}
    \caption{Схема цепи второго порядка}
    \label{fig:circuit_2nd_order}
\end{figure}

\textit{3. Исследование искажений в цепи высокого порядка}
\begin{enumerate}
    \item Соберите схему фильтра высокого порядка согласно указаниям преподавателя. Параметры элементов: $L_1 = 700~\text{мкГн}$, $L_2 = 360~\text{мкГн}$, $C_1 = 4400~\text{пФ}$, $C_2 = 3000~\text{пФ}$, $C_3 = 200~\text{пФ}$, $R_4 = 0,41~\text{кОм}$.
    \item Подайте на вход прямоугольный импульс длительностью $t_{\text{и}} = 10~\text{мкс}$. Снимите осциллограммы напряжений.
    \item Подключите к цепи плоттер Боде (Bode plotter).
    \item Снимите амплитудно-частотную характеристику (АЧХ) в диапазоне частот от $20~\text{кГц}$ до $500~\text{кГц}$.
    \item Занесите полученные значения модуля коэффициента передачи $|H_U(j\omega)|$ для различных частот в \cref{tab:bode_data} протокола.
    \item \textit{Нюанс:} Выбирайте шаг по частоте так, чтобы точно зафиксировать спад характеристики и определить частоту среза на уровне $1/\sqrt{2}$ от максимума.
\end{enumerate}

\begin{figure}[H]
    \centering
    \includegraphics[width=\textwidth]{figs/circuit_high_order.pdf}
    \caption{Схема цепи высокого порядка}
    \label{fig:circuit_high_order}
\end{figure}

% Соответствие изображений из методички:
% fig:circuit_1st_order соответствует Рис.12.1, а (страница 3)
% fig:circuit_2nd_order соответствует Рис.12.1, б (страница 3)
% fig:circuit_high_order соответствует рис.12.2 (схема фильтра на странице 5)
% Полная схема стенда (страница 4) может быть добавлена при необходимости как отдельный рисунок.